% Paste into Overleaf: Conclusions and Future work chapter.
% Keep short; cross-ref Ch5/Ch6; no result dumps; no chapter-role meta.

\chapter{Conclusions and Future work}
\label{ch:conclusionsAndFutureWork}

\section{Conclusions}
\label{sec:conclusions}

This degree project produced a modular research prototype in which structured extraction, multi-turn belief tracking, privacy-restricted aggregate memory, and language-model generation are separated, and in which dialogue decisions are driven by an inspectable intention state.
Relative to the goals in Section~\ref{sec:goals}, the architecture and evaluation apparatus were delivered; the research questions receive a mixed but determinate answer.

For RQ1, explicit intention conditioning did not raise category task match over LLM-only operation on the reported scenarios.
It did increase belief narrowing and clarification under underspecification (Sections~\ref{sec:results-rq1} and~\ref{sec:discussion-intention}).
Robustness here is therefore better understood as controllable uncertainty handling than as higher final category accuracy.

For RQ2, a local aggregate prior can improve match modestly while the prompt privacy audit remains satisfied (Sections~\ref{sec:results-rq2} and~\ref{sec:discussion-memory}).
The privacy side of the claim is supported more strongly than the personalisation side, because the reported memory contrast uses a mild gold-aligned prior rather than priors learned from real booking histories.

Taken together, the work supports a design stance in which the language model realises surface text and the sidecar keeps intention state, clarification policy, and data minimisation enforceable.
It does not support the stronger claim that the modular path must dominate language-model-only decisions on category match.

\section{Limitations}
\label{sec:limitations}

The study is bounded by the delimitations in Section~\ref{sec:delimitations} and by the measurement choices in Chapter~\ref{ch:methods}.
Evaluation uses scripted multi-turn scenarios rather than live Fixmeapp traffic, and category match substitutes for booking completion in the absence of an API or database ground truth.
Emotion enters only as a light dialogue signal, not as a validated affective module.
The Memory-augmented contrast depends on how the prior is initialised; the reported seed is appropriate for ablation, not for claiming deployed personalisation.
Front-end confusions among neighbouring service classes remain a ceiling on scenario-level gains.
No large-scale user study or legal GDPR audit is claimed.

\section{Future work}
\label{sec:futureWork}

Priority follow-ups are those already indicated by the residual errors and by Section~\ref{sec:discussion-improvements}: sharper separation among wellness, fitness, beauty, and medical evidence; clarification questions that target competing high-mass categories; and memory priors estimated from held-out interaction aggregates rather than gold-aligned seeds, under the same prompt allow-list.
A separate track should score generative reply quality without folding it into the sidecar metrics.
Closer to the host setting, limited pilot dialogues with Fixmeapp-style users would test whether the clarify--suggest policy remains acceptable outside scripted packages.

\section{Reflections}
\label{sec:reflections}

The ethical centre of the project is data minimisation in a consumer booking assistant: personal identifiers and detailed preference profiles are kept out of the generative prompt, and personalisation---when used---acts only through a local prior and a compact intention summary.
Making belief and policy explicit also makes failure modes inspectable, which matters when users speak vaguely and when neighbouring service classes share everyday vocabulary.
Economically, the prototype is a design probe for Fixmeapp rather than a production cost study; socially, it aims at clearer booking help under incomplete requests without treating the language model as an unbounded user memory.
Environmental impact is dominated by whatever generative backend is attached; the local stack itself is intentionally lightweight.
These constraints shaped the architecture as much as accuracy targets did, and they should continue to bound any follow-on deployment.
